% Coverage: physical pp. 82--91 (printed pp. 59--68), from the Section 26.2
% heading through the final paragraph on physical p. 91, ending immediately
% before Section 26.3. Equations (26.2.1)--(26.2.41), eleven unnumbered
% displays, citation [1], no footnotes, one centered three-asterisk divider,
% and no figures or tables. Uncertainty: none.

\subsection{General Superfields}

\providecommand{\gradedbracket}[2]{\left[#1,#2\right\}}
\providecommand{\SuperQ}{\mathcal{Q}}

It is straightforward to construct supermultiplets of fields by the direct
technique illustrated in the previous section, but in order to construct
supersymmetric actions we also need to know how to multiply field
supermultiplets to make other supermultiplets. A great deal of work can be
saved by using a formalism invented by Salam and
Strathdee~\hyperref[ch26-ref-1]{[1]}, in which the fields in any
supermultiplet are assembled into a single superfield.

Just as the four-momentum operators $P_\mu$ generate translations of
$x^\mu$, the supersymmetry generators \(Q_\alpha\) and
\(\bar Q_{\dot\alpha}\) may be regarded as generators of translations of
the two independent fermionic superspace coordinates
\(\theta^\alpha\) and \(\bar\theta^{\dot\alpha}\).  These coordinates
anticommute with one another and with fermionic fields, but commute with
$x^\mu$ and all bosonic fields.  The supersymmetry generators have
non-vanishing mixed anticommutators, so they cannot be represented by
fermionic coordinate derivatives alone.  Salam and Strathdee instead found
that their graded brackets with any bosonic or fermionic superfield
$S(x,\theta,\bar\theta)$ may be represented by
\begin{equation}
  \gradedbracket{Q_\alpha}{S}=i\SuperQ_\alpha S,
  \qquad
  \gradedbracket{\bar Q_{\dot\alpha}}{S}
  =i\bar{\SuperQ}_{\dot\alpha}S\,,
  \tag{26.2.1}\label{eq:26.2.1}
\end{equation}
where the superspace differential operators are
\begin{equation}
  \SuperQ_\alpha
  \equiv-\frac{\partial}{\partial\theta^\alpha}
  -i\sigma^\mu_{\alpha\dot\alpha}\bar\theta^{\dot\alpha}\partial_\mu,
  \qquad
  \bar{\SuperQ}_{\dot\alpha}
  \equiv\frac{\partial}{\partial\bar\theta^{\dot\alpha}}
  +i\theta^\alpha\sigma^\mu_{\alpha\dot\alpha}\partial_\mu .
  \tag{26.2.2}\label{eq:26.2.2}
\end{equation}
(All derivatives with respect to fermionic c-number variables are
\emph{left}-derivatives.)  With indices displayed, the first of these
operators is
\begin{equation}
  \SuperQ_\alpha
  =-\partial_\alpha
  -i\sigma^\mu_{\alpha\dot\beta}\bar\theta^{\dot\beta}\partial_\mu .
  \tag{26.2.3}\label{eq:26.2.3}
\end{equation}
Likewise,
\begin{equation}
  \bar{\SuperQ}_{\dot\alpha}
  =\bar\partial_{\dot\alpha}
  +i\theta^\beta\sigma^\mu_{\beta\dot\alpha}\partial_\mu .
  \tag{26.2.4}\label{eq:26.2.4}
\end{equation}
It is straightforward to calculate that
\begin{equation}
  \{\SuperQ_\alpha,\bar{\SuperQ}_{\dot\beta}\}
  =-i\sigma^\mu_{\alpha\dot\beta}\partial_\mu
   -i\sigma^\mu_{\alpha\dot\beta}\partial_\mu .
  \tag{26.2.5}\label{eq:26.2.5}
\end{equation}
The two contributions are equal, and therefore
\begin{equation}
  \{\SuperQ_\alpha,\bar{\SuperQ}_{\dot\beta}\}
  =-2i\sigma^\mu_{\alpha\dot\beta}\partial_\mu .
  \tag{26.2.6}\label{eq:26.2.6}
\end{equation}
Eqs.~\eqref{eq:26.2.6} and \eqref{eq:26.2.1} together with the generalized
Jacobi identities \eqref{eq:25.1.5} show that
\begin{equation}
  [\{Q_\alpha,\bar Q_{\dot\beta}\},S]
  =-\{\SuperQ_\alpha,\bar{\SuperQ}_{\dot\beta}\}S
  =2i\sigma^\mu_{\alpha\dot\beta}\partial_\mu S
  =-2\sigma^\mu_{\alpha\dot\beta}[P_\mu,S]\,,
  \tag{26.2.7}\label{eq:26.2.7}
\end{equation}
in agreement with the anticommutation relation \eqref{eq:25.2.36} in its
two-component form.

It is often more convenient to express the commutation and anticommutation
relations \eqref{eq:26.2.1} as transformation rules.  A transformation
with infinitesimal Weyl parameter \(\xi_\alpha\) changes a superfield by
\begin{equation}
  \delta S=(\xi\SuperQ+\bar\xi\bar{\SuperQ})S
  =-\xi^\alpha\partial_\alpha S
   +\bar\xi_{\dot\alpha}\bar\partial^{\dot\alpha}S
   -i\bigl(\xi\sigma^\mu\bar\theta
     -\theta\sigma^\mu\bar\xi\bigr)\partial_\mu S .
  \tag{26.2.8}\label{eq:26.2.8}
\end{equation}
For example, the left derivatives obey
\begin{equation}
  \partial_\alpha(\theta\theta)=2\theta_\alpha,
  \qquad
  \bar\partial_{\dot\alpha}(\bar\theta\bar\theta)
  =2\bar\theta_{\dot\alpha}.
  \tag{26.2.9}\label{eq:26.2.9}
\end{equation}

The components of \(\theta\) and \(\bar\theta\) anticommute, so the
expansion terminates at order \(\theta^2\bar\theta^2\).  Reducing
Weinberg's original expansion with the comparison dictionary gives
\begin{equation}
  \begin{aligned}
  S(x,\theta,\bar\theta)={}&C
  -i(\theta\omega-\bar\theta\bar\omega)
  -\frac{i}{2}(\theta\theta-\bar\theta\bar\theta)M
  -\frac12(\theta\theta+\bar\theta\bar\theta)N\\
  &+\theta\sigma_\mu\bar\theta\,V^\mu
  -i\theta\theta\,
    \bar\theta\left(\bar\lambda
      -\frac{i}{2}\bar\sigma^\mu\partial_\mu\omega\right)\\
  &+i\bar\theta\bar\theta\,
    \theta\left(\lambda
      -\frac{i}{2}\sigma^\mu\partial_\mu\bar\omega\right)
  +\frac12\theta\theta\bar\theta\bar\theta
    \left(D+\frac12\Box C\right).
  \end{aligned}
  \tag{26.2.10}\label{eq:26.2.10}
\end{equation}
(The derivative terms are separated from \(\lambda\) and \(D\) for later
convenience.) If \(S\) is a scalar, then \(C,M,N,D\) are scalar or
pseudoscalar fields, \(\omega_\alpha,\lambda_\alpha\) are Weyl fields, and
\(V^\mu\) is a vector field.  If \(S=S^\dagger\), the bosonic components
are real and barred spinors are the complex conjugates of the corresponding
unbarred spinors.

For completeness, we spell out the two-component reduction that leads to
the component rules below.  It combines the graded-bracket prescription
\eqref{eq:26.2.1} with the transformation convention
\eqref{eq:26.1.18}; applying \eqref{eq:26.2.1} a second time also checks the
relative sign of the conjugate terms.  The index positions follow
Eq.~\eqref{eq:26.A.3} and the conjugation rule
\eqref{eq:5.4.35}, while the two spacetime-derivative contributions are the
same ones displayed in Eq.~\eqref{eq:26.2.5}. Acting on the expansion
\eqref{eq:26.2.10} gives
\[
  \begin{aligned}
  \delta S
  =(\xi\SuperQ+\bar\xi\bar{\SuperQ})
  \biggl\{&
    C-i(\theta\omega-\bar\theta\bar\omega)
    -\frac{i}{2}(\theta^2-\bar\theta^2)M\\
   &-\frac12(\theta^2+\bar\theta^2)N
    +\theta\sigma^\mu\bar\theta V_\mu+\cdots
  \biggr\}.
  \end{aligned}
\]
The left-derivative convention \eqref{eq:26.2.9} and the first dyad identity
\eqref{eq:26.A.9} give
\[
  \theta_\alpha\theta_\beta
  =\frac12\epsilon_{\alpha\beta}\theta^2,
  \qquad
  \bar\theta_{\dot\alpha}\bar\theta_{\dot\beta}
  =-\frac12\epsilon_{\dot\alpha\dot\beta}\bar\theta^2 .
\]
The same identity \eqref{eq:26.A.9}, contracted once, yields the cubic
reduction \eqref{eq:26.A.16}:
\[
  \theta_\alpha(\theta\chi)
  =-\frac12\theta^2\chi_\alpha,
  \qquad
  \bar\theta_{\dot\alpha}(\bar\theta\bar\chi)
  =-\frac12\bar\theta^2\bar\chi_{\dot\alpha}.
\]
The vector cubic reduction \eqref{eq:26.A.17} is
\[
  (\theta\sigma^\mu\bar\theta)\theta_\alpha
  =-\frac12\theta^2(\sigma^\mu\bar\theta)_\alpha,
  \qquad
  (\theta\sigma^\mu\bar\theta)\bar\theta_{\dot\alpha}
  =-\frac12(\theta\sigma^\mu)_{\dot\alpha}\bar\theta^2 .
\]
At fourth order, Eq.~\eqref{eq:26.A.19} becomes
\[
  \theta_\alpha\bar\theta_{\dot\alpha}
  (\theta\sigma^\mu\bar\theta)
  =\frac14\sigma^\mu_{\alpha\dot\alpha}
    \theta^2\bar\theta^2 .
\]
The exchange rule \eqref{eq:26.A.7} is useful for putting every term in a
common order:
\[
  \xi\sigma^\mu\bar\theta
  =-\bar\theta\bar\sigma^\mu\xi,
  \qquad
  \theta\sigma^\mu\bar\xi
  =-\bar\xi\bar\sigma^\mu\theta .
\]
Collecting the constant and linear terms of \eqref{eq:26.2.10} gives
\[
  \begin{aligned}
  \delta S\big|_{0,1}
  ={}&i(\xi\omega-\bar\xi\bar\omega)\\
   &-i\theta^\alpha
    \left[(-M+iN)\xi_\alpha
      -(\sigma^\mu\bar\xi)_\alpha(\partial_\mu C+iV_\mu)\right]
    +\mathrm{c.c.}
  \end{aligned}
\]
The quadratic part of the same expansion organizes itself as
\[
  \delta S\big|_2
  =-\frac{i}{2}(\theta^2-\bar\theta^2)\delta M
   -\frac12(\theta^2+\bar\theta^2)\delta N
   +\theta\sigma^\mu\bar\theta\,\delta V_\mu .
\]
Finally, the cubic and quartic coefficients retain the shifted combinations
chosen in \eqref{eq:26.2.10}:
\[
  \begin{aligned}
  \delta S\big|_{3,4}={}&
  -i\theta^2\bar\theta\,
   \delta\left(\bar\lambda-\frac{i}{2}
     \bar\sigma^\mu\partial_\mu\omega\right)\\
  &+i\bar\theta^2\theta\,
   \delta\left(\lambda-\frac{i}{2}
     \sigma^\mu\partial_\mu\bar\omega\right)
  +\frac12\theta^2\bar\theta^2
   \delta\left(D+\frac12\Box C\right).
  \end{aligned}
\]

Apply \eqref{eq:26.2.8} to the expansion \eqref{eq:26.2.10}. Reducing the
products with the appendix identities and comparing equal powers of
\(\theta\) and \(\bar\theta\) gives the transformation rules:
\begin{equation}
  \delta C=i(\xi\omega-\bar\xi\bar\omega)\,,
  \tag{26.2.11}\label{eq:26.2.11}
\end{equation}
\begin{equation}
  \delta\omega_\alpha
  =(-M+iN)\xi_\alpha
   -(\sigma^\mu\bar\xi)_\alpha(\partial_\mu C+iV_\mu)\,,
  \tag{26.2.12}\label{eq:26.2.12}
\end{equation}
\begin{equation}
  \delta M=-\xi\left(\lambda-i\sigma^\mu\partial_\mu\bar\omega\right)
  -\bar\xi\left(\bar\lambda
    -i\bar\sigma^\mu\partial_\mu\omega\right)\,,
  \tag{26.2.13}\label{eq:26.2.13}
\end{equation}
\begin{equation}
  \delta N=i\xi\left(\lambda-i\sigma^\mu\partial_\mu\bar\omega\right)
  -i\bar\xi\left(\bar\lambda
    -i\bar\sigma^\mu\partial_\mu\omega\right)\,,
  \tag{26.2.14}\label{eq:26.2.14}
\end{equation}
\begin{equation}
  \delta V_\mu
  =-i\left(\xi\sigma_\mu\bar\lambda
    +\bar\xi\bar\sigma_\mu\lambda\right)
   +\xi\partial_\mu\omega+\bar\xi\partial_\mu\bar\omega\,.
  \tag{26.2.15}\label{eq:26.2.15}
\end{equation}
The cubic and quartic terms, combined with
\eqref{eq:26.2.11} and \eqref{eq:26.2.12}, give the simpler transformation
rules for \(\lambda\) and \(D\):
\begin{equation}
  \delta\lambda_\alpha
  =\left[(\sigma^{\mu\nu})_\alpha{}^\beta
    (\partial_\mu V_\nu-\partial_\nu V_\mu)
    +iD\delta_\alpha{}^\beta\right]\xi_\beta\,,
  \tag{26.2.16}\label{eq:26.2.16}
\end{equation}
\begin{equation}
  \delta D
  =\xi\sigma^\mu\partial_\mu\bar\lambda
   -\bar\xi\bar\sigma^\mu\partial_\mu\lambda\,.
  \tag{26.2.17}\label{eq:26.2.17}
\end{equation}
It is to achieve this simplification that the terms
\(\sigma^\mu\partial_\mu\bar\omega/2\),
\(\bar\sigma^\mu\partial_\mu\omega/2\), and \(\Box C/2\) were separated
from \(\lambda,\bar\lambda\), and \(D\) in the expansion
\eqref{eq:26.2.10}.

The whole point of the superfield formalism is to simplify the task of making
supermultiplets out of other supermultiplets. Given two superfields $S_1$ and
$S_2$ that both satisfy the transformation rule \eqref{eq:26.2.8}, their
product $S\equiv S_1S_2$ satisfies
\begin{equation}
  \begin{aligned}
  \delta S&=-i[\xi Q+\bar\xi\bar Q,S_1S_2]
  =(\delta S_1)S_2+S_1(\delta S_2)\\
  &=\bigl((\xi\SuperQ+\bar\xi\bar{\SuperQ})S_1\bigr)S_2
    +S_1\bigl((\xi\SuperQ+\bar\xi\bar{\SuperQ})S_2\bigr)
  =(\xi\SuperQ+\bar\xi\bar{\SuperQ})S\,,
  \end{aligned}
  \tag{26.2.18}\label{eq:26.2.18}
\end{equation}
and is therefore also a superfield. The calculation uses the exchange and
cubic identities \eqref{eq:26.A.7} and \eqref{eq:26.A.16}. The quartic
terms use Eqs.~\eqref{eq:26.A.18} and \eqref{eq:26.A.19}. The resulting
components are
\begin{equation}
  C=C_1C_2\,,
  \tag{26.2.19}\label{eq:26.2.19}
\end{equation}
\begin{equation}
  \omega=C_1\omega_2+C_2\omega_1\,,
  \tag{26.2.20}\label{eq:26.2.20}
\end{equation}
\begin{equation}
  M=C_1M_2+C_2M_1
  +\frac{i}{2}(\omega_1\omega_2-\bar\omega_1\bar\omega_2)\,,
  \tag{26.2.21}\label{eq:26.2.21}
\end{equation}
\begin{equation}
  N=C_1N_2+C_2N_1
  -\frac12(\omega_1\omega_2+\bar\omega_1\bar\omega_2)\,,
  \tag{26.2.22}\label{eq:26.2.22}
\end{equation}
\begin{equation}
  V^\mu=C_1V_2^\mu+C_2V_1^\mu
  -\frac12\left(
    \omega_1\sigma^\mu\bar\omega_2
    -\bar\omega_1\bar\sigma^\mu\omega_2\right)\,,
  \tag{26.2.23}\label{eq:26.2.23}
\end{equation}
\begin{equation}
  \begin{aligned}
  \lambda_\alpha={}&C_1\lambda_{2\alpha}+C_2\lambda_{1\alpha}
  +\frac{i}{2}(\sigma^\mu\bar\omega_1)_\alpha\partial_\mu C_2
  +\frac{i}{2}(\sigma^\mu\bar\omega_2)_\alpha\partial_\mu C_1\\
  &-\frac12V_{1\mu}(\sigma^\mu\bar\omega_2)_\alpha
  -\frac12V_{2\mu}(\sigma^\mu\bar\omega_1)_\alpha\\
  &+\frac12(N_1-iM_1)\omega_{2\alpha}
  +\frac12(N_2-iM_2)\omega_{1\alpha}\,,
  \end{aligned}
  \tag{26.2.24}\label{eq:26.2.24}
\end{equation}
\begin{equation}
  \begin{aligned}
  D={}&-\partial_\mu C_1\partial^\mu C_2+C_1D_2+C_2D_1
  +M_1M_2+N_1N_2\\
  &-\omega_1\left(\lambda_2
      -\frac{i}{2}\sigma^\mu\partial_\mu\bar\omega_2\right)
   -\bar\omega_1\left(\bar\lambda_2
      -\frac{i}{2}\bar\sigma^\mu\partial_\mu\omega_2\right)\\
  &-\omega_2\left(\lambda_1
      -\frac{i}{2}\sigma^\mu\partial_\mu\bar\omega_1\right)
   -\bar\omega_2\left(\bar\lambda_1
      -\frac{i}{2}\bar\sigma^\mu\partial_\mu\omega_1\right)
  -V_{1\mu}V_2^\mu\,.
  \end{aligned}
  \tag{26.2.25}\label{eq:26.2.25}
\end{equation}
It is trivial that linear combinations of superfields are superfields, in the
same sense, and that spacetime derivatives and complex conjugates of
superfields are superfields. But multiplying a superfield by some function of
$\theta$ or differentiating it with respect to $\theta$ does not in general
yield a superfield. (For instance, $\theta$ itself is evidently not a
superfield, because $\theta$ is a fermionic c-number and therefore commutes
with \(\xi Q+\bar\xi\bar Q\), while $\SuperQ\theta\neq0$.) There is, however, a way of
combining a derivative of a superfield with respect to $\theta$ and
multiplying it by a factor $\theta$ which does yield another superfield.

Consider the covariant superspace derivatives defined by
\begin{equation}
  D_\alpha\equiv-\frac{\partial}{\partial\theta^\alpha}
  +i\sigma^\mu_{\alpha\dot\alpha}\bar\theta^{\dot\alpha}\partial_\mu,
  \qquad
  \bar D_{\dot\alpha}
  \equiv\frac{\partial}{\partial\bar\theta^{\dot\alpha}}
  -i\theta^\alpha\sigma^\mu_{\alpha\dot\alpha}\partial_\mu .
  \tag{26.2.26}\label{eq:26.2.26}
\end{equation}
or more explicitly
\begin{equation}
  D_\alpha=-\partial_\alpha
  +i\sigma^\mu_{\alpha\dot\beta}\bar\theta^{\dot\beta}\partial_\mu,
  \qquad
  \bar D_{\dot\alpha}=\bar\partial_{\dot\alpha}
  -i\theta^\beta\sigma^\mu_{\beta\dot\alpha}\partial_\mu .
  \tag{26.2.27}\label{eq:26.2.27}
\end{equation}
The only difference between the definitions of \(D,\bar D\) and
\(\SuperQ,\bar{\SuperQ}\) is a
change of sign of the terms involving the spacetime derivative. In
consequence of this change of sign, in the mixed anticommutators, instead of
getting two equal terms like those in
Eq.~\eqref{eq:26.2.5} with the same sign, we get them with opposite signs, so
that they cancel:
\begin{equation}
  \{D_\alpha,\SuperQ_\beta\}
  =\{D_\alpha,\bar{\SuperQ}_{\dot\beta}\}
  =\{\bar D_{\dot\alpha},\SuperQ_\beta\}
  =\{\bar D_{\dot\alpha},\bar{\SuperQ}_{\dot\beta}\}=0\,.
  \tag{26.2.28}\label{eq:26.2.28}
\end{equation}
Since \(\xi\) is fermionic, it follows that
\(\xi\SuperQ+\bar\xi\bar{\SuperQ}\) commutes with \(D_\beta\), so if
\(S(x,\theta,\bar\theta)\) is a superfield, then
\begin{equation}
  \delta D_\beta S
  \equiv-i[\xi Q+\bar\xi\bar Q,D_\beta S]
  =D_\beta(\xi\SuperQ+\bar\xi\bar{\SuperQ})S
  =(\xi\SuperQ+\bar\xi\bar{\SuperQ})D_\beta S\,,
  \tag{26.2.29}\label{eq:26.2.29}
\end{equation}
so that \(D_\beta S\) is also a superfield.
\emph{Thus an arbitrary polynomial function of superfields $S$ and their
superderivatives \(D_\beta S,\bar D_{\dot\beta}S\), etc.\
is also a superfield.}

It is not necessary to add here that in constructing a superfield out of other
superfields we can also include their spacetime derivatives, because these can
be obtained from mixed second superderivatives. Their nonzero
anticommutator is
\begin{equation}
  \{D_\alpha,\bar D_{\dot\beta}\}
  =+2i\sigma^\mu_{\alpha\dot\beta}\partial_\mu .
  \tag{26.2.30}\label{eq:26.2.30}
\end{equation}

Now let us consider how to construct a supersymmetric action out of
superfields. There is no such thing as a supersymmetric Lagrangian density,
because the anticommutation relation \eqref{eq:26.2.6} shows that if
$\delta\mathcal{L}=0$ then $\mathcal{L}$ must be a constant. Even if the
Lagrangian density is not supersymmetric, the action will still be
supersymmetric if $\delta\mathcal{L}(x)$ is a derivative, which would not
contribute to $\delta\int\mathcal{L}\,d^4x$. In general, the Lagrangian
density $\mathcal{L}$ can be written as a sum of terms, each of which is some
component of a superfield that is constructed out of elementary superfields
and their superderivatives. Inspection of the transformation rules
\eqref{eq:26.2.11}--\eqref{eq:26.2.17} for the individual components shows
that in the absence of any special conditions on a general superfield, the
only component of such a superfield whose variation is a derivative is the
$D$-component. Also, for the $D$-component of any superfield to be a scalar,
the superfield itself must be a scalar. Therefore, unless there are special
conditions on the individual superfields from which the Lagrangian density is
constructed, a supersymmetric action can only be the integral of the
$D$-term of a scalar superfield $\Lambda$:
\begin{equation}
  I=\int d^4x\,[\Lambda]_D\,.
  \tag{26.2.31}\label{eq:26.2.31}
\end{equation}
But in fact no action of this sort would be physically satisfactory without
special conditions on the superfields from which it is constructed. For a
general superfield \(S(x,\theta,\bar\theta)\), the only sort of
supersymmetric kinematic
action $I_0$ that is bilinear in $S$ and $S^*$ and involves no more than two
derivatives of the component fields is of the form
\begin{equation}
  I_0\mathrel{\propto}\int d^4x\,[S^*S]_D\,.
  \tag{26.2.32}\label{eq:26.2.32}
\end{equation}
From Eq.~\eqref{eq:26.2.25}, we see that $S^*S$ has the $D$-component
\begin{equation}
  \begin{aligned}
  [S^*S]_D={}&-\partial_\mu C^*\partial^\mu C
  +\frac{i}{2}\left(
    \omega\sigma^\mu\partial_\mu\bar\omega
    +\bar\omega\bar\sigma^\mu\partial_\mu\omega\right)\\
  &-\frac{i}{2}\left(
    (\partial_\mu\omega)\sigma^\mu\bar\omega
    +(\partial_\mu\bar\omega)\bar\sigma^\mu\omega\right)\\
  &+C^*D+D^*C
   -(\omega\lambda+\bar\omega\bar\lambda)
   -(\lambda\omega+\bar\lambda\bar\omega)\\
  &+M^*M+N^*N-V_\mu^*V^\mu\,.
  \end{aligned}
  \tag{26.2.33}\label{eq:26.2.33}
\end{equation}
The terms quadratic in $C$ or $\omega$ look promising as the kinematic
Lagrangians for massless fields of spins zero and $1/2$; the final three terms
are harmless; but the terms involving $D$ or $\lambda$ have the disastrous
effect in path integrals of constraining $C$ and $\omega$ to vanish.
Fortunately, as we shall see in the following section, there are constrained
superfields from which we can construct physically sensible actions. The
introduction of these constrained superfields will also open up ways of
constructing supersymmetric terms in the action that are not the
$D$-components of functions of superfields.

If parity is conserved, then the space inversion properties of the component
fields of a superfield will be related by supersymmetry. To work out this
relation, we apply the parity operator $P$ to the commutation/anticommutation
relation \eqref{eq:26.2.1} and use the transformation property
\eqref{eq:25.3.16} of supersymmetry generators, which gives
\begin{equation}
  \begin{aligned}
  -i\gradedbracket{\bar Q_{\dot\alpha}}
    {P^{-1}S(x,\theta,\bar\theta)P}
    &=\SuperQ_\alpha P^{-1}S(x,\theta,\bar\theta)P,\\
  i\gradedbracket{Q_\alpha}
    {P^{-1}S(x,\theta,\bar\theta)P}
    &=\bar{\SuperQ}_{\dot\alpha}
      P^{-1}S(x,\theta,\bar\theta)P .
  \end{aligned}
  \tag{26.2.34}\label{eq:26.2.34}
\end{equation}
Here and below parity identifies the displayed dotted and undotted indices
with \(\sigma^0_{\alpha\dot\alpha}\).
The solution of Eq.~\eqref{eq:26.2.34} for a scalar superfield is of the form
\begin{equation}
  P^{-1}S(x,\theta,\bar\theta)P
  =\eta S(\Lambda_Px,-i\bar\theta,-i\theta)\,,
  \tag{26.2.35}\label{eq:26.2.35}
\end{equation}
Substituting Eq.~\eqref{eq:26.2.35} first into the first relation in
\eqref{eq:26.2.34}, and then into the second relation in
\eqref{eq:26.2.34}, gives
\[
  \begin{aligned}
  \SuperQ_\alpha\,
   S(\Lambda_Px,-i\bar\theta,-i\theta)
  &=-i(\sigma^0)_{\alpha\dot\alpha}
    \bar{\SuperQ}^{\dot\alpha}
    S(\Lambda_Px,-i\bar\theta,-i\theta),\\
  \bar{\SuperQ}_{\dot\alpha}\,
   S(\Lambda_Px,-i\bar\theta,-i\theta)
  &=+i(\sigma^0)_{\alpha\dot\alpha}
    \SuperQ^\alpha
    S(\Lambda_Px,-i\bar\theta,-i\theta).
  \end{aligned}
\]
Thus Eq.~\eqref{eq:26.2.35} solves Eq.~\eqref{eq:26.2.34}; applying
Eq.~\eqref{eq:26.2.35} once more is consistent with the phase and the
index exchange in the displayed parity relation.
Here $\eta$ is some phase (the intrinsic parity of the superfield) and
$\Lambda_Px\equiv(-\mathbf{x},+x^0)$.  Direct substitution into
\eqref{eq:26.2.2} verifies this solution. Using the expansion
\eqref{eq:26.2.10} in
Eq.~\eqref{eq:26.2.35} then gives the space inversion properties of the
component fields:
\begin{equation}
  \begin{aligned}
  P^{-1}C(x)P&=\eta C(\Lambda_Px),\\
  P^{-1}\omega_\alpha(x)P
    &=-i\eta(\sigma^0)_{\alpha\dot\alpha}
      \bar\omega^{\dot\alpha}(\Lambda_Px),\\
  P^{-1}M(x)P&=-\eta M(\Lambda_Px),\\
  P^{-1}N(x)P&=\eta N(\Lambda_Px),\\
  P^{-1}V^\mu(x)P&=-\eta(\Lambda_P)^\mu{}_\nu V^\nu(\Lambda_Px),\\
  P^{-1}\lambda_\alpha(x)P
    &=i\eta(\sigma^0)_{\alpha\dot\alpha}
      \bar\lambda^{\dot\alpha}(\Lambda_Px),\\
  P^{-1}D(x)P&=\eta D(\Lambda_Px).
  \end{aligned}
  \tag{26.2.36}\label{eq:26.2.36}
\end{equation}

\begin{center}
\(\ast\qquad\ast\qquad\ast\)
\end{center}

The general real superfield $S$ involves four real spinless fields $C$, $M$,
$N$, and $D$, plus one real four-vector field $V_\mu$, for a total of eight
independent bosonic field components. For comparison, there are two
Weyl fields \(\omega_\alpha\) and \(\lambda_\alpha\), together with their
conjugates, also for a total of eight independent real field components. The
equality of the numbers of
independent bosonic and fermionic field components holds in general not only
for the unconstrained general superfield studied in this section, but also for
all superfields obtained from the general superfield by imposing
supersymmetric constraints, such as the chiral and other constrained
superfields discussed in the next section.

To see this in general, suppose that we have a representation of the
supersymmetry algebra provided by $N_B$ linearly independent real bosonic
field operators $b_n(x)$ and $N_F$ linearly independent fermionic field
operators $f_k(x)$. We will assume that these fields satisfy only non-trivial
field equations, so that no linear combination of the $b_n$ or $f_k$ with
non-vanishing coefficients satisfies a homogeneous linear field equation.
Consider a Hermitian supersymmetry generator \(Q(u)\), defined as
\begin{equation}
  Q(u)\equiv u^\alpha Q_\alpha
    +\bar u_{\dot\alpha}\bar Q^{\dot\alpha}\,,
  \tag{26.2.37}\label{eq:26.2.37}
\end{equation}
where \(u_\alpha\) is an ordinary commuting numerical Weyl spinor
(\emph{not} an anticommuting c-number) and
\(\bar u_{\dot\alpha}=u_\alpha^*\). (For extended supersymmetry, we could
use any one of the \(Q_{r\alpha},\bar Q_{r\dot\alpha}\) pairs.) In order for
the $b_n$ and $f_k$ to furnish a representation of the supersymmetry algebra,
we must have
\begin{equation}
  [Q(u),b_n]=i\sum_k q_{nk}(\partial)f_k\,,
  \tag{26.2.38}\label{eq:26.2.38}
\end{equation}
\begin{equation}
  \{Q(u),f_k\}=\sum_n p_{kn}(\partial)b_n\,,
  \tag{26.2.39}\label{eq:26.2.39}
\end{equation}
for some matrix differential operators $q(\partial)$ and $p(\partial)$.
Taking the anticommutator of Eq.~\eqref{eq:26.2.38} and the commutator of
Eq.~\eqref{eq:26.2.39} with $Q(u)$ gives
\begin{equation}
  [Q^2(u),b_n]
  =i\sum_m\bigl(q(\partial)p(\partial)\bigr)_{nm}b_m\,,
  \tag{26.2.40}\label{eq:26.2.40}
\end{equation}
\begin{equation}
  [Q^2(u),f_k]
  =i\sum_\ell\bigl(p(\partial)q(\partial)\bigr)_{k\ell}f_\ell\,.
  \tag{26.2.41}\label{eq:26.2.41}
\end{equation}
The anticommutation relation \eqref{eq:25.2.36} or \eqref{eq:25.2.38} gives
the square of \(Q(u)\) as
\(Q^2(u)=-2(u\sigma^\mu\bar u)P_\mu\). Hence the
square matrices $p(\partial)q(\partial)$ and $q(\partial)p(\partial)$ must
both be non-singular, because if there were any non-vanishing coefficients
$c_n(\partial)$ or $d_k(\partial)$ for which
$\sum_n c_n(\partial)(q(\partial)p(\partial))_{nm}=0$ or
$\sum_k d_k(\partial)(p(\partial)q(\partial))_{k\ell}=0$ then $b_n$ or $f_k$
would satisfy the homogeneous linear field equations
\[
  (u\sigma^\mu\bar u)\partial_\mu
    \sum_n c_n(\partial)b_n=0
  \qquad\text{or}\qquad
  (u\sigma^\mu\bar u)\partial_\mu
    \sum_k d_k(\partial)f_k=0\,,
\]
in contradiction with our assumption that the fields do not satisfy such
field equations. In order for $qp$ to be non-singular we must have
$N_F\geq N_B$, and in order for $pq$ to be non-singular we must have
$N_B\geq N_F$, so we can conclude that $N_B=N_F$. Also the square matrices
$q$ and $p$ must both be non-singular, so the complex conjugate of
Eq.~\eqref{eq:26.2.38} tells us that $f^*=q^{*-1}qf$, so that the number of
independent fermion fields is $N_F$ rather than $2N_F$, and is therefore
equal to the number $N_B$ of independent boson fields, as was to be shown.
